\documentclass{exam}
\usepackage{../../commonheader}

\discnumber{10}
\title{Efficiency and Caching}
\date{October 26, 2026 -- October 30, 2026}

% Exercises and lessons are imported from the topic bank when possible.
\begin{document}
\maketitle

\begin{meta}
\small
\textbf{Material Scope and Older Exam Notation}
\begin{itemize}
    \item \textbf{Efficiency scope:} Big O and orders of growth were not taught this semester. Focus on memoization and list-building choices; no growth classifications or formal notation are required.
    \item \textbf{Out of scope:} \lstinline{__repr__}, linked list mutation, and tree mutation.
    \item \textbf{Still in scope:} \lstinline{__str__} and everything else.
    \item Older exam problems use older notation. Translate \lstinline{t.is_leaf()} to \lstinline{is_leaf(t)}, and \lstinline{Link.empty} to \lstinline{()}.
\end{itemize}

\textbf{Example Timeline}
\begin{itemize}
    \item Check-in and caching overview [10 min]. Trace a cache miss and a hit. Locate the dictionary and explain why it survives between calls.
    \item Q1: repeated Fibonacci calls [8 min]. Draw the calls and identify repeated inputs.
    \item Q2: implement and trace a cache [17 min]. Compare an initial request, a repeated request, and the next Fibonacci number.
    \item Q3: list-building choices [10 min]. Explain copying, shifting, and output order.
    \item Wrap-up [5 min]. Ask students what repeated work each technique avoids.
\end{itemize}

\textbf{Teaching Goals}
\begin{itemize}
    \item Tree-recursive Fibonacci repeats the same calls many times. Memoization saves each result in a dictionary keyed by the argument; a repeated request becomes a lookup.
    \item Every recursive call must check the shared cache. Each Fibonacci value is computed only once, although it may be requested many times. Q2 starts with the base-case results already stored.
    \item Keep the dictionary outside the returned function so it persists across calls. Two calls to \lstinline{make_fib} create separate caches.
    \item Test membership, not truthiness: 0 is a valid cached result. Stored results use memory, and reuse is appropriate only while the inputs still determine the same result.
    \item Appending and list comprehensions efficiently build lists. Repeated concatenation copies the growing list; front insertion shifts existing elements. Appending occasionally resizes storage, but does not copy the entire list on every append.
    \item A comprehension still evaluates its expression for every element. Discuss concrete work avoided, not runtime formulas or growth classifications.
\end{itemize}
\textbf{Optional Midterm Review: Problem Selection}
These three topic-bank problems have not appeared on earlier Fall 2026 worksheets. Efficiency remains this week's main topic; choose review problems as time permits rather than adding all three to the timeline above.
\begin{itemize}
    \item Q4: \lstinline{Poll} and \lstinline{tally} [12--15 min]. The bank credits Fall 2018 Midterm 2 Q5(a). Students coordinate a shared class list, per-instance dictionaries, and a filtered comprehension. This is comparable to the object-state reasoning in Fall 2023 Midterm 2 Q3 (Library/Book), without requiring out-of-scope representation methods.
    \item Q5: \lstinline{max_subset_sum} [10--12 min]. Students track both the remaining list and the number of elements still needed, then combine include/exclude branches with \lstinline{max}. Compare Fall 2023 Midterm 2 Q4(a), \lstinline{fit}: both require coordinated recursive arguments and base cases. This bank problem is somewhat more scaffolded.
    \item Q6: \lstinline{SkipMachine} [8--10 min]. Track each generator's local \lstinline{current}, its instance's \lstinline{skip}, and the shared \lstinline{SkipMachine.skip}. Fall 2023 Midterm 2 Q1 similarly tests lazy execution and iterator consumption; this problem adds shared class state. Calling \lstinline{generate} does not execute its body. Each \lstinline{next}
    resumes after the previous \lstinline{yield}; evaluate the parts in order.
\end{itemize}
Difficulty comparisons are teaching judgments, not official exam ratings. Reference: \href{https://website.cs61a.org/exam/fa23/mt2/61a-fa23-mt2_sol.pdf}{Fall 2023 Midterm 2 solutions}.
\end{meta}

\section{Efficiency}
\begingroup
\subimport{../../topics/growth/text/}{caching-overview.tex}
\endgroup
\begin{questions}
    \subimport{../../topics/growth/caching/}{fib-repeated-work.tex}

    \subimport{../../topics/growth/caching/}{memoized-fib.tex}

    \subimport{../../topics/growth/caching/}{list-building.tex}
\end{questions}

\newpage
\section{Optional Midterm Review}
\begin{questions}
    \setcounter{question}{3}
    \subimport{../../topics/oop/hard/}{midterm-elections.tex}

    \newpage
    \subimport{../../topics/recursion/medium/}{max-subset-sum.tex}

    \newpage
    \subimport{../../topics/iterators/easy/}{skipmachine.tex}
\end{questions}

\end{document}
